\documentclass[12pt]{article}
\usepackage{polyglossia}
\usepackage{fontspec}
\usepackage[export]{adjustbox}
\usepackage{amsfonts, amsmath, amssymb}
\usepackage{annotate-equations}
\usepackage{bm}
\usepackage{bookmark}
% \usepackage{breqn}
\usepackage{derivative}
\usepackage{empheq}
\usepackage[inline]{enumitem}
\usepackage{fancyhdr}
\usepackage{geometry}
\usepackage{graphicx}
\usepackage{kantlipsum}
\usepackage{lastpage}
\usepackage{linebreaker} % line-breaker algorithm in LuaLaTeX
\usepackage{lualatex-math} % Fixes for mathematics
\usepackage{mathtools}
\usepackage{microtype}
\usepackage{mismath}
\usepackage{nicefrac}
\usepackage[hyperref]{ntheorem}
% \usepackage{phfqit} % Bra-ket notation
% \usepackage{physics}
\usepackage[final]{showlabels}
\usepackage{simples-matrices}
\usepackage{siunitx}
\usepackage{soul}
\usepackage{spbmark}
\usepackage{titling}
\usepackage{tabularray}
\usepackage{tasks}
\usepackage{tensor}
\usepackage{xcolor}
\usepackage{xfrac}
\usepackage{xparse}
\usepackage{xspace}
\usepackage{hyperref}
% Polyglossia config
\setdefaultlanguage[variant=mexican]{spanish}
\setotherlanguage{english}
\usepackage{cleveref}

%%%%%%%%%%%%%%%%%%%%%%%%%%%%%%%%%%%%%%%%%%%%%%%%%%%%%%%%%%
%%%%%%%%%%%%%%%%%%%%%%%%%%%%%%%%%%%%%%%%%%%%%%%%%%%%%%%%%%
%%%%%%%%%%%%%%%%%%Configuraciones extras%%%%%%%%%%%%%%%%%%
\definecolor{base3}{RGB}{253, 246, 227}%
\definecolor{pinkwave}{RGB}{255, 0, 128}%
\definecolor{customBlue}{RGB}{11, 61, 98}%
\pagecolor{base3}
\graphicspath{{img/}}
\setlength{\parindent}{2em} % Sangría
\setlength{\parskip}{0.5em} % Espacio entre párrafos
\linespread{1.1} % line spacing
\setlength{\jot}{10pt} % Space between lines in multiline eqs
% \crefname{equation}{ec.}{ecs.} % Equation's cross-reference name
\crefname{equation}{}{} % Equation's cross-reference name

% Line-breaker config
\linebreakersetup {
    maxtolerance = 90,
    maxemergencystretch = 1em,
    maxcycles = 4
}

%%%%%%%%%%%%%%%%%%Theorem environments%%%%%%%%%%%%%%%%%%
% Configuración de ambiente para problema
\theoremstyle{break}
\theoremheaderfont{\Large\normalfont\bfseries}
\theorembodyfont{\normalfont}
\theoremseparator{\bigskip} % Spacing between header and body
\theorempreskip{1.5em}
\theorempostskip{\topsep\bigskip}
\theorempostwork{
    \color{customBlue} \hrule width \hsize height 2pt \kern 1mm \hrule width \hsize
    }
\newtheorem{exercise}{Problema}
\counterwithin{equation}{exercise}

\makeatletter
\LetLtxMacro\old@exercise\exercise
\renewcommand*{\exercise}{%
  \@ifnextchar[{%
    \exercise@opt
  }{%
    \old@exercise
    \exercise@bookmark{}%
  }%
}
\def\exercise@opt[#1]{%
  \old@exercise[{#1}]%
  \exercise@bookmark{: #1}%
}
\newcommand*{\exercise@bookmark}[1]{%
  \bookmark[dest=\@currentHref]{Problema \theexercise#1}%
}
\makeatother

% Configuración de ambiente para solución
\theoremstyle{nonumberbreak}
\theoremheaderfont{\Large\normalfont\bfseries}
\theorembodyfont{\normalfont}
\theoremseparator{\medskip}
\theorempreskip{1em}
\theorempostskip{\topsep\medskip}
\newtheorem{solution}{Solución}

%%%%%%%%%%%%%%%%%%%%%%%%%%%%%%%%%%%%%%%%%%%%%%%%%%%%%%%%


% Configuración del paquete hyperref
\hypersetup{
    colorlinks = true,%
    linkcolor={[rgb]{0,0.2,0.6}},%
    citecolor={[rgb]{0,0.6,0.2}},%
    filecolor={[rgb]{0.8,0,0.8}},%
    urlcolor={[rgb]{0.8,0,0.8}},%
    runcolor={[rgb]{0.8,0,0.8}},% 
    menucolor={[rgb]{0,0.2,0.6}},%
    linkbordercolor={[rgb]{0,0.2,0.6}},%
    citebordercolor={[rgb]{0,0.6,0.2}},%
    filebordercolor={[rgb]{0.8,0,0.8}},%
    urlbordercolor={[rgb]{0.8,0,0.8}},%
    runbordercolor={[rgb]{0.8,0,0.8}},%
    menubordercolor={[rgb]{0,0.2,0.6}},% 
    pdftitle={Tarea 1},%
    pdfauthor={López Merino Marcos},%
    pdfsubject={Mecánica Analítica},%
    pdfkeywords={Facultad de Ciencias, UNAM, Mecánica Analítica},%
    unicode = true%
}

%%%%%%%%%%%%%%%%%%siunitx configuration%%%%%%%%%%%%%%%%%%
% Configuración del paquete siunitx
\sisetup{
	output-decimal-marker = {.}, 
	per-mode = symbol-or-fraction,
	separate-uncertainty = false,
	exponent-product = \cross,
    % inter-unit-product = \ensuremath{{}\vdot{}}
}

% Declaring new units
\DeclareSIUnit\kilogram{\kilo\gram}

%%%%%%%%%%%%%%%%%%%%%%%%%%%%%%%%%%%%%%%%%%%%%%%%%%%%%%%%

% Geometría del documento
\geometry{
    letterpaper,
    top = 0.6in,
    bottom = 0.8in,
    left = 0.6in,
    right = 0.6in,
    footskip = 38pt
}

%%%%%%%%%%%%%%%%%%Nuevos comandos%%%%%%%%%%%%%%%%%%
\newcommand*{\group}{8195}
\newcommand*{\classname}{Mecánica Analítica}
\newcommand*{\homeworknumber}{Tarea 1}
\newcommand*{\name}{Marcos López Merino}

% unit vector i
\newcommand*{\uveci}{{\bm{\hat{\textnormal{\bfseries\i}}}}}
% unit vector j
\newcommand*{\uvecj}{{\bm{\hat{\textnormal{\bfseries\j}}}}}
% unit vector
\DeclareRobustCommand{\uvec}[1]{{%
  \ifcsname uvec#1\endcsname
     \csname uvec#1\endcsname
   \else
    \bm{\hat{\mathbf{#1}}}%
   \fi
}}% 
\newcommand{\idest}{\emph{i.e.},\xspace} % id est
% Espacio vectorial, e.g., ℝ, ℂ, ℕ, etc.
\NewDocumentCommand{\vecspace}{m o}{%
  \IfValueTF{#2}{%
    \mathbb{#1}^{#2}%
  }{%
    \mathbb{#1}%
  }%
}

% unit vector k
\newcommand*{\uvectk}{\uvec{k}}

\newcommand*{\observer}{\mathcal{O}}
\newcommand*{\primeobserver}{\mathcal{O}^{\prime}}
\newcommand*{\biprimeobserver}{\mathcal{O}^{\prime\prime}}
\newcommand*{\triprimeobserver}{\mathcal{O}^{\prime\prime\prime}}
\newcommand*{\inlinesol}{\vspace*{10pt}\textbf{Solución}\vspace*{10pt}}
\renewcommand*{\grad}{\nabla}
\renewcommand*{\div}{\nabla \cdot}
\renewcommand*{\curl}{\nabla \mul}
\newcommand*{\laplacian}{\nabla^{2}}
\newcommand*{\qq}{\quad\quad}
\DeclareDifferential{\cdif}{\grad}[style-notation=single]
\newcommand*{\dotprod}[2]{#1 \mathbin{\cdot} #2\xspace}
\newcommand*{\crossprod}[2]{#1 \mathbin{\mul} #2\xspace}
\newcommand*{\basevect}[1]{\uvec{e}_{#1}\xspace}

% \newcommand*{\change}[1]{\Delta{#1}\xspace}

\NewDocumentCommand{\change}{ s o m }{%
    \IfBooleanTF{#1}{%
        \Delta {#3^{\prime}}
    }{%
        \Delta {#3}
    }
}
\NewDocumentCommand{\interval}{ s o m }{%
    \IfBooleanTF{#1}{%
        (\Delta {#3^{\prime}})^{2}
    }{%
        (\Delta {#3})^{2}
    }
}
\DeclarePairedDelimiterX\set[1]\lbrace\rbrace{#1}
% \newcommand{\crefrangeconjunction}{--}
% \newcommand{\crefpairconjunction}{\xspace y\xspace}

% Defining a variant of Aboxed command from mathtools
\makeatletter
\newcommand*\Acolorboxed[2][pinkwave]{%
   \let\bgroup{\romannumeral-`}%
   \@Acolorboxed{#1}#2&&\ENDDNE
}
\def\@Acolorboxed#1#2&#3&#4\ENDDNE{%
  \ifnum0=`{}\fi
  \setbox\z@\hbox{$\displaystyle#2{}\m@th$\kern\fboxsep \kern\fboxrule}%
  \edef\@tempa{\kern\wd\z@ & \kern-\the\wd\z@ \fboxsep\the\fboxsep \fboxrule\the\fboxrule}%
  \@tempa
  \fcolorbox{#1}{base3}{\m@th$\displaystyle#2#3$}%
} 
\makeatother

%%%%%%%%%%%%%%%%%%Portada y configuración%%%%%%%%%%%%%%%%%%
% Configuración de portada
\setlength{\droptitle}{-60pt} % raise the title

% Portada
\title{
    \textbf{\homeworknumber}\\
    \normalsize\vspace{0.1in}\small{\textbf{Entrega}:~\today}
    \vspace{-1.5in}
}
\author{}
\date{}

%%%%%%%%%%%%%%%%%%Header and footer%%%%%%%%%%%%%%%%%%
\setlength{\headheight}{15.2pt}
\pagestyle{fancy}
\lhead{Grupo \group, Sem. 2023-1}
\chead{\classname}
\rhead{\name}
\lfoot{\includegraphics[scale = 0.2, valign = c]{LogoFCUNAMcolor.pdf}}
% \lfoot{\includegraphics[scale = 0.1, valign = c]{example-image}}
\cfoot{\homeworknumber}
\rfoot{Pág. \thepage \hspace{1pt} de \pageref{LastPage}}

\renewcommand{\headrulewidth}{0.5pt}
\renewcommand{\footrulewidth}{0.5pt}
%%%%%%%%%%%%%%%%%%%%%%%%%%%%%%%%%%%%%%%%%%%%%%%%%%%%%%%%%%
%%%%%%%%%%%%%%%%%%%%%%%%%%%%%%%%%%%%%%%%%%%%%%%%%%%%%%%%%%

\begin{document}
    \maketitle
    \thispagestyle{fancy}

    % Problema 1
    \begin{exercise}[\textcolor{red}{1.5}]
      Comprobar la ecuación

      \begin{equation*}
        \crossprod{\vect{A}}{\left(\crossprod{\vect{B}}{\vect{C}}\right)} = \left(\dotprod{\vect{A}}{\vect{C}}\right)\vect{B} - \left(\dotprod{\vect{A}}{\vect{B}}\right)\vect{C}.
      \end{equation*}

      \textbf{desarrollando ambos miembros} en función de las componentes cartesianas de los vectores. Hállese el triple producto vectorial \(\vect{A} \mul \left(\vect{B} \mul \vect{C}\right)\) de los tres vectores

      \begin{equation*}
        \vect{A} = 3\uveci + \uvecj + 5\uvectk,\quad \vect{B} = -\uveci + 4\uvecj - 2\uvectk,\quad \vect{C} = 2\uveci - \uvecj + \uvectk.
      \end{equation*}
    \end{exercise}
    
    % Problema 2
    \begin{exercise}[\textcolor{red}{1.11}]
      Hállese el conjunto recíproco de vectores del conjunto de vectores no coplanares.

      \begin{enumerate}[label = \alph*)]
        \item 
        
          \begin{align*}
            \vect{b}_{1} &= 2\uveci + \uvecj + 4\uvectk,\\
            \vect{b}_{2} &= \uveci + 3\uvectk,\\
            \vect{b}_{3} &= -3\uveci - 4\uvecj - \uvectk.
          \end{align*}

        \item 
          
          \begin{align*}
            \vect{b}_{1} &= \uveci + 2\uvecj + 2\uvectk,\\
            \vect{b}_{2} &= -3\uveci - 4\uvecj + 2\uvectk,\\
            \vect{b}_{3} &= 2\uveci - \uvecj - \uvectk.
          \end{align*}
      \end{enumerate}
    \end{exercise}
    
    % Problema 3
    \begin{exercise}[\textcolor{red}{1.12}]
      \begin{enumerate}[label = \roman*)]
        \item Exprese los vectores
        
        \begin{equation*}
          \vect{A} = 2\uveci - 4\uvecj + 3\uvectk,\quad \vect{B} = -3\uveci + 2\uvecj - \uvectk,
        \end{equation*}

        en función de una suma lineal de los vectores no coplanares del problema anterior (\textcolor{red}{1.11}) y en función de una suma lineal de sus vectores recíprocos.

        \item Evaluar \(\dotprod{\vect{A}}{\vect{B}}\) y \(\crossprod{\vect{A}}{\vect{B}}\) utilizando los productos escalares \(\alpha_{i},\ \beta_{i},\ \alpha_{i}^{*}\) y \(\beta_{i}^{*}\) de \(\vect{A}\) y \(\vect{B}\) con el conjunto de vectores no coplanares encontrados en i) y sus vectores recíprocos. Compárese la respuesta con los productos escalar y vectorial de estos tres vectores evaluados en función de las componentes cartesianas de \(\vect{A}\) y \(\vect{B}\).
        
      \end{enumerate}
    \end{exercise}
    
    % Problema 4
    \begin{exercise}[\textcolor{red}{1.16}]
      La derivada del vector \(\vect{A}(t)\) de magnitud constante se demostró que puede ser expresada como un producto vectorial
      
      \begin{equation*}
        \odv{\vect{A}(t)}{t} = \crossprod{\vect{\omega}}{\vect{A}(t)},
      \end{equation*}

      Sin embargo, \(\vect{\omega}\) no es única, ya que la suma del término \(c\vect{A}\) a \(\vect{\omega}\) producirá el mismo resultado. Por otro lado, la derivada de dos vectores de magnitud constante puede determinar una \(\vect{\omega}\) única, en función de la cual sus derivadas son expresables en la forma

      \begin{equation*}
        \odv*{\vect{A}(t)}{t} = \crossprod{\vect{\omega}}{\vect{A}}\qquad \text{y} \qquad \odv*{\vect{B}(t)}{t} = \crossprod{\vect{\omega}}{\vect{B}}.
      \end{equation*}

      Considerando los vectores unidad

      \begin{align*}
        \basevect{1} &= \sin\alpha t\cos\beta t\uveci + \sin\alpha t\sin\beta t\uvecj + \cos\alpha t\uvectk,\\
        \basevect{2} &= \cos\alpha t\cos\beta t\uveci + \cos\alpha t\sin\beta t\uvecj - \sin\alpha t\uvectk,\\
        \basevect{3} &= -\sin\beta t\uveci + \cos\beta t\uvecj,
      \end{align*}

      hállese el vector de la velocidad angular \(\vect{\omega}\) que satisface las ecuaciones

      \begin{equation*}
        \odv*{\basevect{1}}{t} = \crossprod{\vect{\omega}}{\basevect{1}}\qquad \text{y} \qquad \odv*{\basevect{2}}{t} = \crossprod{\vect{\omega}}{\basevect{2}}.
      \end{equation*}

      Demostrar que la misma \(\vect{\omega}\) también nos da

      \begin{equation*}
        \odv*{\basevect{3}}{t} = \crossprod{\vect{\omega}}{\basevect{3}}.
      \end{equation*}
    \end{exercise}
    
    % Problema 5
    \begin{exercise}[\textcolor{red}{1.18}]
      \begin{enumerate}[label = \alph*)]
        \item Demostrar que los productos escalares \(\alpha_{i}\) y \(\alpha_{i}^{*}\) del vector \(\vect{A}\) por los vectores base \(\vect{b}_{i}\) y sus recíprocos \(\vect{\bm{b}}_{i}\), están relacionados linealmente como se indica por
        
        \begin{equation*}
          \alpha_{i}^{*} = \sum_{j}g_{ij}^{*}\alpha_{j}\qquad \text{y}\qquad \alpha_{i} = \sum_{j}g_{ij}\alpha_{j}^{*},
        \end{equation*}

        Exprésense los escalares \(g_{ij}\) y \(g_{ij}^{*}\) en función de los vectores base \(\vect{b}_{i}\) y los vectores recíprocos \(\vect{\bm{b}}_{i}\). (Obsérvese que los \(g_{ij}\) son las componentes covariantes del tensor métrico y las \(g_{ij}^{*}\) son las componentes contravariantes del mismo.)

        \item Demostrar que el producto escalar de dos vectores \(\vect{A}\) y \(\vect{B}\) como
        
        \begin{equation*}
          \dotprod{\vect{A}}{\vect{B}} = \sum_{ij}g_{ij}^{*}\alpha_{i}\beta_{j} = \sum_{ij}g_{ij}\alpha_{i}^{*}\beta_{j}^{*}.
        \end{equation*}
      \end{enumerate}
    \end{exercise}
    
    % Problema 6
    \begin{exercise}[\textcolor{red}{1.20}]
      Comprobar para las ecuaciones

      \begin{equation*}
        \vect{b}_{1} = \dfrac{\crossprod{\vect{\bm{b}}_{2}}{\vect{\bm{b}}_{3}}}{\dotprod{\vect{\bm{b}_{1}}}{\crossprod{\vect{\bm{b}}_{2}}{\vect{\bm{b}}_{3}}}},\quad
        \vect{b}_{2} = \dfrac{\crossprod{\vect{\bm{b}}_{3}}{\vect{\bm{b}}_{1}}}{\dotprod{\vect{\bm{b}_{1}}}{\crossprod{\vect{\bm{b}}_{2}}{\vect{\bm{b}}_{3}}}},
        \quad
        \vect{b}_{3} = \dfrac{\crossprod{\vect{\bm{b}}_{1}}{\vect{\bm{b}}_{2}}}{\dotprod{\vect{\bm{b}_{1}}}{\crossprod{\vect{\bm{b}}_{2}}{\vect{\bm{b}}_{3}}}}
      \end{equation*}

      que los vectores \(\vect{b}_{1},\ \vect{b}_{2}\) y \(\vect{b}_{3}\) son recíprocos \(\vect{\bm{b}}_{1},\ \vect{\bm{b}}_{2}\) y \(\vect{\bm{b}}_{3}\).
    \end{exercise}
    
    % Problema 7
    \begin{exercise}[\textcolor{red}{2.2}]
      Obtener las expresiones de las componentes polares de la velocidad y la aceleración de las partículas cuyo vector bidimensional de posición es:

      \begin{equation*}
        r = \dfrac{5}{2 - \cos\phi},\qquad \phi = \omega t.
      \end{equation*}
    \end{exercise}
    
    % Problema 8
    \begin{exercise}[\textcolor{red}{2.3}]
      Hállese, explícitamente, la velocidad y la aceleración de una partícula que se mueve a lo largo de una hélice circular definida por las coordenadas cilíndricas.

      \begin{equation*}
        \rho = a,\qquad \phi = \omega t\qquad \text{y} \qquad z = -ct,\qquad a,\ b,\ c \in \vecspace{R}.
      \end{equation*}
    \end{exercise}
\end{document}